\section{The Night of the Black Cows}

There are two forms of darkness. One is where we encounter God. The other is the night in which all cows are black. That is, the night during which no one is able to make intelligent distinctions. The former is considered obscurantism by our educated contemporaries. The latter is considered ``enlightenment" by the same people.

\paragraph{Divine Darkness}
Shortly after I posted \textit{The Philosophy of the Future}\footnote{\url{https://gornahoor.net/?p=7379}} last week, I received the journal of \textit{Death to the World}\footnote{\url{http://www.deathtotheworld.com/}}. Through synchronicity, it contained a passage from Saint Gregory of Nyssa, titled \emph{Divine Darkness}\footnote{\url{https://gornahoor.net/?p=6071}}, on the prophet \textbf{Moses}. He wrote:

\begin{quotex}
What does it mean that Moses entered the darkness and then saw God in it? What is now recounted seems somehow to be contradictory to the first theophany, for then that the divine was beheld in light, but now He is seen in darkness. 

\end{quotex}
Although religious knowledge is first received as light, he continues:

\begin{quotex}
as the mind progresses and, through an ever greater and more perfect diligence, comes to apprehend reality, as it approaches more nearly to contemplation, it sees more clearly what of the divine nature is uncontemplated. 

\end{quotex}
Gregory explains the method of ``depth", not to mention the trial to give up one's false beliefs:

\begin{quotex}
For leaving behind everything that is observed, not only what sense comprehends, but also what the intelligence thinks it sees, it keeps on penetrating deeper until by the intelligence's yearning for understanding it gains access to the invisible and the incomprehensible, and there it sees God. This is the true knowledge of what is thought; this is the seeing that consists in not seeing because that which is sought transcends all knowledge, being separated on all sides by incomprehensibility, as by a kind of darkness. 

\end{quotex}
\paragraph{Consensus Reality}
\textbf{Jacques Ellul} explained that the purpose of mass education is to make it possible for propaganda to be absorbed more effectively. Rather than the method of depth which requires the abandonment of all that one falsely observes and thinks, the method of propaganda is just the opposite. It seeks to fill the mind with unjustified and unjustifiable opinions. This requires a certain level of sophistication so that is why the intelligent and the degreed are nearly always in accord with modern consensus reality.

The under-educated and the less intelligent, on the other hand, are unable to attain that level, so are confused by it, should it ever cross their mind. In that case, they often gravitate to marginal, and often ridiculous, conspiracy theories. This does not challenge the modern worldview, but rather reinforces it, since the propagandized can gloat over their intelligence compared to the uneducated.

If we let X represent consensus reality and we define a set of mental assumptions that imply X, we can represent it like this:

$[A,B,C,…] \Rightarrow  X$, where $A$ is atheism, for example, and B, C, etc are other assumptions, whether explicitly acknowledged or implicitly believed.

The remarkable thing is that we find that certain assumptions, which seem to be foundational to one's worldview, make absolutely no difference. For example, Mary says she believes X because of her Christian principles, while the atheist professor Waldo believes the same thing. If the above formula refers to the professor, we can express Mary's world view in this one:

$[\bar A,B,C,…] \Rightarrow  X$, where $\bar A$ is not-$A$

Then logically, assumption A is irrelevant to coming to conclusion X. Yet Mary and Waldo insist the opposite. If all the cows are black, then A, B, and C are irrelevant. Everything looks like X.

\paragraph{Gnosis}
The super-correct seem to think that gnosis is a threat to their religion. I'll quote \textbf{Frithjof Schuon} to help define the concept:

\begin{quotex}
a key idea of Christianity is that it means ``God became what we are, in order that we might become what He is" (\textbf{St Irenaeus}); Heaven became earth that earth might become Heaven. Christ re-enacts in the outward historical world what is being enacted at all times in the inner world of the soul. In man the Spirit becomes the ego in order that the ego may become pure Spirit; the Spirit becomes ego by incarnating in the mind in the form of intellection, of truth, and the ego becomes the Spirit through uniting with it. 

\end{quotex}
There is clearly no threat here; the esoteric way is orthogonal to the exoteric way. Since the exoteric path is sufficient for salvation, it simply cannot see the point of esoterism. Hence, it can only be believed as irrelevant at best, or heretical at worst. Their eyes cannot penetrate the darkness.

What happens historically, is also experienced interiorly. As above, so below. As the inside, so the outside.

\paragraph{The Kaleidsocopic Bliss of Heaven}
Unlike humans, angels do not share a common nature; each one has his own nature, utterly distinct from every other angel. Moreover, they have no individual form. Hence, our way of experiencing them is different. Schuon explains it very nicely:

\begin{quotex}
The angels are created of light, of a supra-formal substance; the differences between angels are like those between colors, sounds, or perfumes, not those between forms … 

\end{quotex}
That is what Heaven is like. It is a multi-dimensional kaleidoscope with constantly changing colors, sound, scents, and experiences undreamed of. As the Depth is Infinite, there can be no end to the chants of the heavenly hosts, and no boredom. Only more and more profundity toward the Divine Darkness where God is hiding in plain sight.



\flrightit{Posted on 2014-07-04 by Cologero }

\begin{center}* * *\end{center}

\begin{footnotesize}\begin{sffamily}



\texttt{Wayne Ferguson on 2014-07-05 at 21:37 said: }

Apropos of ``The Kaleidsocopic Bliss of Heaven", what do you make of this?

\url{http://www.youtube.com/watch?v=7agl-sNLXMI}


\hfill

\texttt{Cologero on 2014-07-08 at 09:16 said: }

Wayne, I don't know if you are addressing that question at large or to me, since I wrote precisely what I make of this. The metaphysics, as described, is sound. The vision is personal and has no authority beyond that. But it can be compared to the rash of books out today about boys and men who claim to have visited ``heaven".


\hfill

\texttt{Wayne Ferguson on 2014-07-08 at 16:06 said: }

Thank you for the reply, Cologero–I was wondering what you (or anyone) might make of the YouTube clip of the housewife on LSD in the 1956 experiment (I shared the link in my first comment). Her description– and evident joy –in that clip, sounds a lot like the ``Kaleidsocopic Bliss of Heaven" in your post. While I don't want to disparage anyone's personal experience or traditional faith/hope, I don't put much stock at all in booksellers like Eben Alexander, et al. On the other hand, your description– and that of the housewife in the YouTube clip –reminds me of William Blake, who I do respect:

``If the doors of perception were cleansed every thing would appear to man as it is, Infinite. For man has closed himself up, till he sees all things thro' narrow chinks of his cavern." 

Thanks again…


\hfill

\texttt{Cologero on 2014-07-10 at 23:57 said: }

Curious, Wayne. I suppose her experience is the material analog of authentic spiritual experience. That is, the experience is passive, caused from the ``outside", without requiring any fundamental change in the woman. I recall in the 60s in Boston, when Leary and Ram Dass were at Harvard and LSD was not yet illegal, that there was a debate on public TV between Leary and some biologist. I don't recall what the latter said. Nevertheless, at most LSD can only indicate certain possibilities, and is no substitute for the real thing.


\hfill

\texttt{Max on 2014-07-12 at 12:10 said: }

Drugs may function to break away from certain illusions, they may also cause worse illusions. In a previous age and through a regular tradition the hallucinogenic experience could probably be guided by a superior in the correct direction. This is not possible anymore or very unlikely. Shamans degenerated into solely dealing with demons and spirits, forgetting the All-Father. The risks increase with the amount of psychic residue as the age proceeds. Hallucinogens probably bring far more pitfalls than opportunities and should be stayed away from, which is the view of Charles Upton who has written about the topic in his books which I recommend on the subject. 

There is reason to believe that those substances also contributed to opening and expanding ``the fissures in the great wall" that Guenon writes of. Psychic experiments coincide in time with accelerating technological and nuclear developments: space exploration, computers, internet and so on which likely are connected in a manner most are ignorant of. As all cows are black for our educated contemporaries as you point out, they never ask themselves where new technology really comes from or what had to be offered for it. There is no such thing as something being discovered or invented just by chance at a specific moment in time.


\end{sffamily}\end{footnotesize}
